    \begin{equation}
      \tag{1.247}
      \begin{split}
        \left(1 - \frac{ip\Delta x'}{\hbar}\right)
        &=
        \int dx'
        \mathscr{J}(\Delta x')\ket{x'}\braket{x'}{\alpha}
        \\
        &=
        \int dx'
        \ket{x' + \Delta x'}\braket{x'}{\alpha}
        \\
        &=
        \int dx'
        \ket{x'}\braket{x' -  \Delta x'}{\alpha}
        \\
        &=
        \int dx'
        \ket{x'}
        \left(
        \braket{x'}{\alpha} - \Delta x' \pdv{}{x'}\braket{x'}{\alpha}
        \right)
        \\
      \end{split}
    \end{equation}

    What he does there at the end is important.  That last
    substitution is the result of taking a first-order Taylor series
    which essentially boils down to

    \begin{equation}
      f(x + dx) = f(x) + dx\pdv{}{x}f(x)
    \end{equation}

    Assuming that the function is differentiable, then the principle
    of local linearity holds and the first-order Taylor series is
    rigorous and complete.  This version, of course, does \emph{not}
    hold up under macroscopic translation, so this only works for
    complete differentials.
